% =====================================================================
%  Universidad Central del Ecuador  -  Sistemas Colaborativos
%  Simulacion Multiagente - Zona de Cobro por Congestion (La Carolina)
%  Plantilla a dos columnas, limpia y profesional.
%  Compilar con latexmk (recomendado) o manualmente para resolver
%  referencias cruzadas y la bibliografia (BibTeX, ver references.bib):
%     latexmk -pdf main.tex
%     o bien: pdflatex main && bibtex main && pdflatex main && pdflatex main
%  Requiere la carpeta images/ junto a este archivo con las 9 figuras:
%     fig1_mapa_poligono, fig02_entorno_gama_eb, fig3_flujo_vehicular,
%     fig4_velocidad_media, fig5_decisiones_bdi, fig6_modal_shift,
%     fig7_benchmark_londres, fig8_equidad_recaudacion, fig9_tarifa_por_tipo
%     (mas el logoIEEEAccess.png, usado en el encabezado de cada pagina)
% =====================================================================
\documentclass[twocolumn,10pt]{article}

% ----------------------- Codificacion e idioma -----------------------
\usepackage[utf8]{inputenc}
\usepackage[T1]{fontenc}
\usepackage{lmodern}
\usepackage[spanish,es-tabla]{babel}   % "Cuadro" -> "Tabla" automatico

% ----------------------------- Matematicas ---------------------------
\usepackage{amsmath,amssymb,amsfonts}

% ------------------------ Graficos y tablas --------------------------
\usepackage{graphicx}
\usepackage{booktabs}      % reglas de tabla de calidad cientifica
\usepackage{array}
\usepackage{tabularx}      % columnas de ancho automatico (X)
\usepackage{float}
\usepackage{placeins}      % \FloatBarrier: evita que las figuras se pasen a Referencias
\usepackage{caption}
\usepackage{pifont}        % simbolos de verificacion (\ding)

% --------------------------- Diseno y color --------------------------
\usepackage{xcolor}
\usepackage{titlesec}
\usepackage{microtype}     % evita desbordes de texto en columnas
\usepackage{enumitem}
\usepackage[a4paper,total={175mm,245mm},left=18mm,top=20mm]{geometry}
\usepackage{fancyhdr}
\usepackage[hidelinks]{hyperref}

\setlength{\columnsep}{6mm}

% --- Paleta institucional (azul academico) ---
\definecolor{azulUCE}{RGB}{31,97,164}
\definecolor{azulOscuro}{RGB}{18,52,86}
\definecolor{grisSuave}{RGB}{244,247,250}
\definecolor{verdeOK}{RGB}{22,140,74}
\definecolor{rojoNo}{RGB}{192,57,43}

\newcommand{\yes}{\textcolor{verdeOK}{\ding{51}}}
\newcommand{\no}{\textcolor{rojoNo}{\ding{55}}}

\hypersetup{colorlinks=true,linkcolor=azulOscuro,citecolor=azulUCE,urlcolor=azulUCE}

% --- Que las leyendas digan "Fig." y "Tabla" (estilo IEEE, en espanol) ---
\AtBeginDocument{%
  \renewcommand{\figurename}{Fig.}\renewcommand{\tablename}{Tabla}%
  \ifdefined\decimalpoint\decimalpoint\fi% punto decimal tambien en modo matematico (consistencia con el texto)
}

% --- Estilo de secciones ---
\titleformat{\section}{\normalsize\bfseries\color{azulUCE}}{\thesection.}{0.5em}{}
\titleformat{\subsection}{\small\bfseries\color{azulOscuro}}{\thesubsection}{0.5em}{}
\titlespacing*{\section}{0pt}{11pt}{5pt}
\titlespacing*{\subsection}{0pt}{8pt}{3pt}

% --- Estilo de leyendas (figuras/tablas) ---
\captionsetup{font=small,labelfont={bf,color=azulUCE},labelsep=period}

% --- Manejo de flotantes (muchas figuras/tablas a doble columna) ---
\setcounter{topnumber}{2}
\setcounter{dbltopnumber}{2}
\setcounter{totalnumber}{4}
\renewcommand{\topfraction}{0.92}
\renewcommand{\dbltopfraction}{0.92}
\renewcommand{\bottomfraction}{0.8}
\renewcommand{\textfraction}{0.06}
\renewcommand{\floatpagefraction}{0.7}
\renewcommand{\dblfloatpagefraction}{0.7}

% --- Encabezado tipo IEEE: logo arriba a la derecha + linea negra, en cada pagina ---
\pagestyle{fancy}
\fancyhf{}
\fancyhead[R]{\includegraphics[height=1.1cm]{images/logoIEEEAccess.png}}
\renewcommand{\headrulewidth}{1pt}
\fancyfoot[C]{\small\thepage}
\setlength{\headheight}{34pt}
\setlength{\headsep}{10pt}

% --- Cintillo institucional (texto de facultad, solo portada) ---
\newcommand{\membrete}{%
  \begin{center}
    {\large\scshape\bfseries Universidad Central del Ecuador}\\[2pt]
    {\small Facultad de Ingenier\'ia y Ciencias Aplicadas \,\textbar\, Carrera de Computaci\'on}\\[2pt]
    {\small Asignatura: Sistemas Colaborativos (TGP09BFT03)}%
  \end{center}
}

\begin{document}

% =====================================================================
%  CABECERA, TITULO, AUTORES Y RESUMEN  (ancho completo)
% =====================================================================
\twocolumn[%
  \membrete
  \vspace{14pt}
  \begin{center}
    {\LARGE\bfseries\color{azulOscuro}
     Simulaci\'on Multiagente del Impacto de una Zona de Cobro por
     Congesti\'on en el Sector Parque La Carolina:\\[3pt]
     Evaluaci\'on de Pol\'iticas de Movilidad Urbana en Quito Mediante
     Arquitectura BDI\par}

    \vspace{12pt}
    {\large Ariel Elizalde, Edison Enr\'iquez, Angelo Silva,
            Stalin Acurio, Alexis Troya\par}
    \vspace{5pt}
    {\small Carrera de Computaci\'on, Facultad de Ingenier\'ia y Ciencias Aplicadas\\
            Universidad Central del Ecuador, Quito, Ecuador\\[2pt]
            \texttt{\{amelizalde, erenriquezp, amsilvac1, svacurio, aetroyab\}@uce.edu.ec}\par}
    \vspace{4pt}
    {\small\itshape Per\'iodo Acad\'emico 2026 --- Noveno Semestre\par}

    \vspace{14pt}
    % ---------- Resumen a ancho completo ----------
    \begin{minipage}{0.93\textwidth}
      {\color{azulUCE}\rule{\linewidth}{0.5pt}}\par\vspace{6pt}
      {\small\noindent\textbf{Resumen---}En este trabajo se desarroll\'o y evalu\'o un
      modelo de Simulaci\'on Multiagente (SMA) orientado a estimar el impacto de una
      zona de cobro por congesti\'on en el pol\'igono comercial y financiero del Parque
      La Carolina, en el Distrito Metropolitano de Quito. Sobre la plataforma GAMA y
      con datos geoespaciales reales procesados en QGIS, se integr\'o la red vial
      vectorial y se modelaron conductores heterog\'eneos ---diferenciados por nivel
      socioecon\'omico y por cinco tipos de veh\'iculo--- parametrizados con reglas
      deliberativas Creencia--Deseo--Intenci\'on (BDI), cuyas funciones de utilidad
      multicriterio ponderaron tiempo de viaje, costo financiero y comodidad. Se
      contrast\'o un escenario base sin peaje (E0) frente a un esquema de peaje por
      franja horaria con tarifas diferenciadas por tipo de veh\'iculo y ajuste din\'amico
      mediante un agente gestor (EB), tomando como referente el London Congestion
      Charge. Los resultados mostraron que el peaje redujo el flujo vehicular interno en
      hora pico un 17.2\,\% e increment\'o la transferencia modal al Metro un 45\,\% en esa
      franja (de 9.8\,\% a 14.3\,\%), \'unico efecto estad\'isticamente significativo
      ($p=0.004$); el agente gestor ejerci\'o su rango din\'amico, escalando la tarifa
      base del auto de \$2.00 al techo de \$3.00 ante densidad alta y recaudando
      \$12\,194 en el d\'ia. Sin embargo, la reducci\'on no alcanz\'o el umbral del 20\,\%
      planteado y se observ\'o un desplazamiento perif\'erico del 52.7\,\%, por lo que la
      hip\'otesis principal no se confirm\'o. Se concluy\'o que el instrumento oper\'o
      correctamente, con beneficios moderados pero efectos colaterales relevantes bajo
      la tarifa evaluada.\par}
      \vspace{6pt}
      {\small\noindent\textbf{Palabras clave---}Simulaci\'on multiagente, arquitectura
      BDI, tarificaci\'on por congesti\'on, GAMA, QGIS, modelado basado en agentes,
      movilidad urbana, Quito.\par}
      \vspace{6pt}
      {\color{azulUCE}\rule{\linewidth}{0.5pt}}
    \end{minipage}
    \vspace{16pt}
  \end{center}
]

% =====================================================================
%  1. INTRODUCCION   (tiempo presente)
% =====================================================================
\section{Introducci\'on}
El Distrito Metropolitano de Quito (DMQ) registra una saturaci\'on cr\'itica de su
infraestructura vial por el crecimiento sostenido de un parque automotor que supera
las 600\,000 unidades en circulaci\'on activa \cite{inec2023}. Esta presi\'on genera
externalidades negativas que degradan la productividad de la capital: p\'erdida de
horas-persona en congesti\'on, mayor consumo energ\'etico y aumento de emisiones. El
sector del Parque La Carolina ---delimitado por las avenidas Naciones Unidas, Amazonas,
De los Shyris, Eloy Alfaro y 6 de Diciembre--- concentra el n\'ucleo financiero de la
urbe y, por la convergencia de viajes pendulares, carece de mecanismos automatizados de
gesti\'on adaptativa de la demanda.

Ante la propuesta conceptual del Municipio del DMQ de implementar una zona de
tarificaci\'on vial en este sector, se identifica un vac\'io de evidencia computacional
publicada que eval\'ue cuantitativamente la viabilidad de la medida en el contexto
espec\'ifico de Quito. Los modelos agregados tradicionales fallan al promediar las
conductas individuales; por ello, el Modelado Basado en Agentes (ABM) resulta id\'oneo,
pues representa la heterogeneidad de los conductores, captura la emergencia de
fen\'omenos sist\'emicos ---como el desplazamiento del tr\'afico hacia v\'ias
perif\'ericas--- y permite evaluar impactos de equidad socioecon\'omica
\cite{gama2023,etrr2023}. La tarificaci\'on din\'amica mediante agentes deliberativos es,
adem\'as, una l\'inea metodol\'ogica activa en la ingenier\'ia de transporte sostenible
\cite{magtToll2024,kaddoura2019}.

El \textbf{objetivo general} es desarrollar un SMA con arquitectura BDI e integraci\'on
de datos geoespaciales reales que simule y eval\'ue el impacto de una zona de cobro por
congesti\'on en el sector Parque La Carolina. Como \textbf{objetivos espec\'ificos} se
plantean: (i)~integrar la red vial real del pol\'igono en un entorno de simulaci\'on
geolocalizado; (ii)~modelar el comportamiento deliberativo de distintos perfiles de
conductores mediante funciones de utilidad multicriterio; y (iii)~contrastar
cuantitativamente un escenario base frente a un escenario de peaje din\'amico,
compar\'andolo con el referente emp\'irico de Londres.

El estudio se gu\'ia por la \textbf{hip\'otesis} de que la implementaci\'on de una tarifa
din\'amica reduce el volumen vehicular de paso interno en al menos un 20\,\%, incrementa
el uso intermodal del Metro de Quito y no provoca un colapso por desv\'io en las arterias
perif\'ericas. El umbral del 20\,\% se establece tomando como referencia el $-27$\,\%
observado en el London Congestion Charge durante sus primeros seis meses (2003)
\cite{london2020,london2024}. El aporte de este trabajo es ofrecer la primera evidencia
computacional de tarificaci\'on por congesti\'on, basada en agentes BDI con datos
geoespaciales reales, para una ciudad andina latinoamericana con un sistema de Metro de
operaci\'on reciente.

% =====================================================================
%  2. MARCO TEORICO   (tiempo presente)
% =====================================================================
\section{Marco Te\'orico y Trabajo Relacionado}
El control de la congesti\'on urbana mediante sistemas multiagente constituye una
l\'inea activa de la ingenier\'ia de transporte sostenible \cite{gama2023,etrr2023}. Los
enfoques reactivos tradicionales son superados por arquitecturas deliberativas
cognitivas como el modelo BDI (Creencias, Deseos, Intenciones), formalizado por Rao y
Georgeff \cite{raoGeorgeff1995}. En un ecosistema de transporte urbano:

\begin{itemize}[leftmargin=1.3em,itemsep=2pt,topsep=2pt]
  \item \textbf{Creencias} --- estado informacional del conductor: posici\'on en la red,
        tarifa vigente en los puntos de control y restricciones activas por el sistema
        local de ``tercera placa''.
  \item \textbf{Deseos} --- objetivos individuales: minimizar el tiempo de viaje, reducir
        los costos financieros directos y resguardar el confort.
  \item \textbf{Intenciones} --- el plan comprometido tras la deliberaci\'on: mantener la
        ruta pagando el peaje, desviarse (\textit{rerouting}) por la periferia o ejecutar
        un cambio intermodal hacia el Metro.
\end{itemize}

Alineado con los fundamentos de los sistemas colaborativos y la ingenier\'ia de sistemas
complejos \cite{gomez2013}, el entorno de simulaci\'on se clasifica como parcialmente
observable, din\'amico y continuo. La coordinaci\'on sist\'emica se alcanza mediante un
agente facilitador (Gestor AMT) que emite se\~nales tarifarias en funci\'on de la
densidad global del \'area, buscando optimizar la eficiencia de la red de forma
distribuida, en l\'inea con los esquemas de \textit{congestion pricing} basados en
agentes reportados en la literatura \cite{magtToll2024,kaddoura2019}.

% =====================================================================
%  3. METODOLOGIA   (tiempo pasado, orden cronologico)
% =====================================================================
\section{Metodolog\'ia}
La construcci\'on y calibraci\'on del simulador se ejecut\'o de forma cronol\'ogica en tres
fases: preparaci\'on geoespacial, especificaci\'on conductual BDI y dise\~no experimental
de escenarios.

\subsection{Pipeline de integraci\'on geoespacial}
Se delimit\'o el pol\'igono de estudio del Parque La Carolina en QGIS~3.x. Mediante el
complemento QuickOSM se descarg\'o la red vial vectorial real desde OpenStreetMap y se
aplic\'o una limpieza topol\'ogica para eliminar nodos duplicados, corregir el sentido de
los flujos y estructurar el grafo de conectividad. Se vectorizaron las capas de los cinco
accesos semaforizados (puntos de control C1--C5) y la ubicaci\'on de las estaciones del
Metro adyacentes (I\~naquito y La Carolina). La base geoespacial se export\'o en formato
Shapefile bajo la proyecci\'on UTM Zona 17S (EPSG:32717), asegurando precisi\'on m\'etrica
al importarse en GAMA mediante la capa \texttt{road\_network}. El entorno geom\'etrico se
detalla en la Fig.~\ref{fig:mapa}.

\begin{figure}[!t]
  \centering
  \includegraphics[width=\columnwidth]{images/fig1_mapa_poligono.png}
  \caption{Pol\'igono georreferenciado del Parque La Carolina (QGIS): red vial vectorial
           real de la zona de cobro por congesti\'on, sobre la cartograf\'ia base de Quito.}
  \label{fig:mapa}
\end{figure}

\subsection{Parametrizaci\'on y funci\'on de utilidad de los agentes}
Se codificaron en lenguaje GAML conductores BDI gobernados por dos identidades
superpuestas. La primera es el \textbf{nivel socioecon\'omico (NSE)} ---Alto, Medio y
Bajo---, que fija la disposici\'on a pagar ($wtp$), los pesos de la funci\'on de utilidad y
el umbral de congesti\'on que dispara la deliberaci\'on (Tabla~\ref{tab:nse}). La segunda
es el \textbf{tipo de veh\'iculo}, que determina la velocidad libre, el factor de
ocupaci\'on de v\'ia y la exoneraci\'on o no del peaje (Tabla~\ref{tab:vehiculos}). Esta
doble parametrizaci\'on es el mecanismo que habilita el an\'alisis de equidad: distintos
perfiles enfrentan el mismo cobro con distinta elasticidad.

% ---------- TABLA I: parametros NSE (ancho completo) ----------
\begin{table*}[!t]
  \centering
  \caption{Par\'ametros por nivel socioecon\'omico (NSE).}
  \label{tab:nse}
  \small
  \renewcommand{\arraystretch}{1.25}
  \begin{tabular}{@{}lcccccc@{}}
    \toprule
    \textbf{NSE} & \textbf{Proporci\'on} & \textbf{$wtp$ (USD)} & \textbf{$w_1$ (tiempo)} &
    \textbf{$w_2$ (costo)} & \textbf{$w_3$ (comodidad)} & \textbf{Umbral congesti\'on}\\
    \midrule
    Alto  & 15\,\% & 3.00 & 0.35 & 0.15 & 0.50 & 0.40\\
    Medio & 45\,\% & 2.25 & 0.35 & 0.35 & 0.30 & 0.55\\
    Bajo  & 40\,\% & 0.50 & 0.20 & 0.65 & 0.15 & 0.70\\
    \bottomrule
  \end{tabular}
\end{table*}

% ---------- TABLA II: tipos de vehiculo (ancho completo) ----------
\begin{table*}[!t]
  \centering
  \caption{Tipos de veh\'iculo, flota inicial (300 agentes) y tratamiento tarifario.}
  \label{tab:vehiculos}
  \small
  \renewcommand{\arraystretch}{1.3}
  \begin{tabularx}{\textwidth}{@{}lcccX@{}}
    \toprule
    \textbf{Tipo} & \textbf{Flota} & \textbf{Factor de v\'ia} &
    \textbf{Tarifa pico (USD)} & \textbf{Observaci\'on}\\
    \midrule
    Moto  & 45 (15\,\%)  & 0.3 & \textbf{Exonerado} & Mayor agilidad; umbral reducido.\\
    Auto  & 165 (55\,\%) & 1.0 & 2.00 & Tarifa base de referencia.\\
    SUV   & 45 (15\,\%)  & 1.5 & 3.00 & Paga 1.5$\times$ la tarifa base del auto.\\
    Bus   & 30 (10\,\%)  & 3.0 & \textbf{Exonerado} & Ruta fija; NSE Bajo forzado.\\
    Carga & 15 (5\,\%)   & 2.5 & 3.00 & Mayor impacto vial (1.5$\times$ el auto).\\
    \bottomrule
  \end{tabularx}
\end{table*}

La toma de decisiones intermodal y de enrutamiento se gobern\'o mediante la funci\'on de
utilidad multicriterio de la ecuaci\'on~\eqref{eq:utilidad}:

\begin{equation}
U(\text{alt}) = w_1\,\frac{1}{T_v} + w_2\,\frac{1}{C} + w_3\,C_p
\label{eq:utilidad}
\end{equation}

\noindent donde $T_v$ es el tiempo estimado de viaje, $C$ el costo econ\'omico directo
asociado a la tarifa y $C_p$ el coeficiente de comodidad percibida del modo. El conductor
selecciona la intenci\'on de m\'axima utilidad entre RUTA\_DIRECTA, REROUTEAR y METRO;
cuando la tarifa efectiva supera su disposici\'on a pagar ($\tau_{ef} > wtp$), la utilidad
de la ruta directa se penaliza, forzando una alternativa.

A diferencia de la propuesta conceptual inicial de una tarifa plana, el modelo final
aplic\'o una \textbf{se\~nal tarifaria diferenciada por tipo de veh\'iculo y activa solo en
franja pico}, seg\'un la ecuaci\'on~\eqref{eq:tarifa}:

\begin{equation}
\tau_k(t)=
\begin{cases}
\$2.00 & k=\text{Auto},\; t\in\text{pico}\\[2pt]
\$3.00 & k\in\{\text{SUV, Carga}\},\; t\in\text{pico}\\[2pt]
\$0.00 & k\in\{\text{Moto, Bus}\}\;\lor\; t\notin\text{pico}
\end{cases}
\label{eq:tarifa}
\end{equation}

\noindent con franjas pico de 07:00--10:00 y 17:00--20:00. La tarifa que cada agente
realmente paga (\textbf{tarifa efectiva}) escala la tarifa base de su tipo por un factor
din\'amico $f_g$ fijado por el gestor, seg\'un la ecuaci\'on~\eqref{eq:efectiva}:

\begin{equation}
\tau_{ef}=\tau_k\cdot f_g,\qquad f_g=\frac{\tau_{\text{din}}}{\tau_{\text{Auto}}}
\label{eq:efectiva}
\end{equation}

\noindent donde $\tau_{\text{din}}$ es la tarifa de auto vigente que define el gestor y
$\tau_{ef}=0$ si el tipo est\'a exonerado o $t\notin$ pico. As\'i, una elevaci\'on del
factor $f_g$ encarece proporcionalmente a todos los tipos no exonerados (p.\,ej., con
$f_g=1.5$ el auto paga \$3.00 y SUV/Carga \$4.50). Sobre esta se\~nal base oper\'o un
\textbf{agente gestor (GestorAMT)} con arquitectura BDI deliberativa que, cada 30 ciclos,
ajust\'o din\'amicamente $\tau_{\text{din}}$ en el rango $[\$0.50,\ \$3.00]$ en pasos de
\$0.25, seleccionando entre las intenciones MANTENER, SUBIR, BAJAR o SUSPENDER en
funci\'on de la densidad del pol\'igono, la velocidad media y la saturaci\'on del Metro, y
difundi\'endola a los puntos de control mediante un patr\'on FIPA.

El indicador de desempe\~no central se defini\'o como la reducci\'on porcentual del flujo
vehicular interno entre el escenario base (0) y el regulatorio (B), expresada en la
ecuaci\'on~\eqref{eq:reduccion}:

\begin{equation}
\Delta Q_{int}=\frac{Q_{int}^{(0)}-Q_{int}^{(B)}}{Q_{int}^{(0)}}\times 100\,\%
\label{eq:reduccion}
\end{equation}

Para evaluar la equidad, se calcul\'o el \textbf{\'indice de Gini modal} sobre la
participaci\'on del Metro por estrato NSE, seg\'un la ecuaci\'on~\eqref{eq:gini}:

\begin{equation}
\begin{aligned}
G   &= \frac{n+1-2\sum_{i=1}^{n} L_i}{n},\\[2pt]
L_i &= \frac{\sum_{j=1}^{i} x_{(j)}}{\sum_{j=1}^{n} x_{(j)}}
\end{aligned}
\label{eq:gini}
\end{equation}

\noindent donde $x_{(j)}$ son las cuotas de uso del Metro ordenadas por estrato. Para este
an\'alisis se aplic\'o una correcci\'on que descuenta la contribuci\'on de los buses
(siempre RUTA\_DIRECTA y NSE Bajo por dise\~no), evitando inflar artificialmente ese
contador.

\begin{figure}[!t]
  \centering
  \includegraphics[width=\columnwidth]{images/fig02_entorno_gama_eb.png}
  \caption{Entorno de GAMA Platform ejecutando el escenario EB al cierre del d\'ia
           (22:00): flota heterog\'enea sobre la red de La Carolina, puntos de control
           C1--C11, zona de cobro activa y los paneles de equidad NSE y recaudaci\'on
           acumulada (\$12\,208).}
  \label{fig:gama}
\end{figure}

\subsection{Especificaci\'on de escenarios}
Se ejecutaron computacionalmente dos escenarios en un horizonte simulado de 06:00 a 22:00
(paso de 10\,s, exportaci\'on de m\'etricas cada $\sim$15\,min sim):
\begin{enumerate}[leftmargin=1.4em,itemsep=2pt,topsep=2pt]
  \item \textbf{Escenario E0 (l\'inea base):} situaci\'on vial actual del DMQ con esquema
        de Pico y Placa activo, sin cobros por circulaci\'on interna y con la oferta base
        del Metro. Los puntos de control existen pero no cobran.
  \item \textbf{Escenario EB (peaje por franja horaria):} misma flota y red que E0, m\'as
        la maquinaria de cobro. El peaje se activ\'o solo en franja pico con tarifas
        diferenciadas por tipo de veh\'iculo (Auto \$2.00, SUV y Carga \$3.00; Moto y Bus
        exonerados), gestionado din\'amicamente por el GestorAMT en el rango
        $[\$0.50,\ \$3.00]$. Oper\'o con tercera placa activa ($\approx$20\,\% de la flota,
        restricci\'on 06:00--20:00) y \textit{modal shift} hacia dos estaciones del Metro
        (I\~naquito y La Carolina).
\end{enumerate}

Ambos escenarios partieron de una flota de 300 conductores que se increment\'o mediante
inyecciones de tr\'afico en hora pico (hasta $\sim$1928 agentes activos al cierre del
d\'ia), reproduciendo las oleadas de ingreso al cord\'on. La velocidad media se midi\'o
dentro de la zona de cobro como $V_{libre}\times$ coeficiente de congesti\'on, para no
diluir el efecto del peaje con la red perif\'erica. Cada escenario produjo 64 registros
temporales que la canalizaci\'on en Python (pandas/NumPy/SciPy) proces\'o para generar las
tablas comparativas, las pruebas estad\'isticas y las figuras.

% =====================================================================
%  4. RESULTADOS
% =====================================================================
\section{Resultados}
Las series temporales y m\'etricas agregadas se exportaron en archivos tabulares tras la
ejecuci\'on de los ciclos en GAMA. Los resultados consolidados se presentan en la
Tabla~\ref{tab:metricas}.

% ---------- TABLA 1: metricas E0 vs EB (ancho completo) ----------
\begin{table*}[!t]
  \centering
  \caption{Comparativa de m\'etricas de rendimiento de la red vial (E0 vs EB).}
  \label{tab:metricas}
  \small
  \renewcommand{\arraystretch}{1.25}
  \begin{tabularx}{\textwidth}{@{}Xrrrc@{}}
    \toprule
    \textbf{Variable evaluada} & \textbf{E0 (base)} & \textbf{EB (peaje)} &
    \textbf{$\Delta$\,\%} & \textbf{\textquestiondown Mejora?}\\
    \midrule
    Flujo interno --- total (veh/h)                       & 3913.00 & 3839.00 & $-1.89$   & \yes\\
    \textbf{Flujo interno --- hora pico (veh/h)}          & 3635.69 & 3011.85 & $\mathbf{-17.16}$ & \yes\\
    Desplazamiento perif\'erico (veh/h)                   & 344.50  & 525.94  & $+52.67$  & \no\\
    Velocidad media --- total (km/h)                      & 30.99   & 31.14   & $+0.48$   & \yes\\
    \textbf{Velocidad media --- hora pico (km/h)}         & 31.35   & 32.08   & $\mathbf{+2.34}$  & \yes\\
    Modal shift $\rightarrow$ Metro --- total (\%)        & 9.45    & 6.80    & $-28.05$  & \no\\
    \textbf{Modal shift $\rightarrow$ Metro --- pico (\%)}& 9.83    & 14.29   & $\mathbf{+45.34}$ & \yes\\
    $\Delta$ Gini modal --- hora pico (0--1)              & 0.12    & 0.10    & $-16.19$  & \yes\\
    $\Delta$ Gini modal --- promedio diario (0--1)        & 0.15    & 0.34    & $+125.16$ & \no\\
    Recaudaci\'on estimada (USD/h)                        & 0.00    & 762.15  & ---       & ---\\
    \bottomrule
  \end{tabularx}
\end{table*}

La respuesta din\'amica del volumen interno ante la activaci\'on de la ventana tarifaria
se grafica en la Fig.~\ref{fig:flujo}, y la velocidad media comparada en la
Fig.~\ref{fig:velocidad}. La bifurcaci\'on de intenciones (\textit{modal shift}) entre
perfiles se visualiza en las Figs.~\ref{fig:bdi} y~\ref{fig:modal}, el contraste con el
referente de Londres en la Fig.~\ref{fig:benchmark}, y la equidad, recaudaci\'on y
diferenciaci\'on tarifaria en las Figs.~\ref{fig:equidad} y~\ref{fig:tarifatipo}.

% ---------- FIGURAS 3 y 4: flujo y velocidad (doble columna) ----------
\begin{figure*}[!t]
  \centering
  \begin{minipage}{0.47\textwidth}
    \centering
    \includegraphics[width=\linewidth]{images/fig3_flujo_vehicular.png}
    \caption{Curvas temporales del volumen de veh\'iculos activos dentro del per\'imetro de
             La Carolina (E0 vs EB) a lo largo del d\'ia.}
    \label{fig:flujo}
  \end{minipage}\hfill
  \begin{minipage}{0.47\textwidth}
    \centering
    \includegraphics[width=\linewidth]{images/fig4_velocidad_media.png}
    \caption{Evoluci\'on de la velocidad media dentro del pol\'igono de cobro (E0 vs EB).}
    \label{fig:velocidad}
  \end{minipage}
\end{figure*}

% ---------- FIGURAS 5 y 6: decisiones BDI y modal shift (doble columna) ----------
\begin{figure*}[!t]
  \centering
  \begin{minipage}{0.47\textwidth}
    \centering
    \includegraphics[width=\linewidth]{images/fig5_decisiones_bdi.png}
    \caption{Distribuci\'on de intenciones de la flota (ruta directa / \textit{rerouting}
             / Metro) por escenario.}
    \label{fig:bdi}
  \end{minipage}\hfill
  \begin{minipage}{0.47\textwidth}
    \centering
    \includegraphics[width=\linewidth]{images/fig6_modal_shift.png}
    \caption{Transferencia modal hacia el Metro por franja horaria (E0 vs EB).}
    \label{fig:modal}
  \end{minipage}
\end{figure*}

Durante la hora pico, el peaje redujo el flujo interno de 3635.7 a 3011.9\,veh/h
($-17.16$\,\%) y elev\'o la velocidad media de 31.35 a 32.08\,km/h ($+2.34$\,\%). La
transferencia modal al Metro en pico creci\'o de 9.83\,\% a 14.29\,\% ($+45.34$\,\%). En
contraste, el desplazamiento perif\'erico aument\'o un 52.67\,\% (de 344.5 a
525.9\,veh/h), evidenciando la reasignaci\'on de parte del tr\'afico hacia v\'ias no
cobradas. El \'indice de Gini modal en hora pico se mantuvo pr\'acticamente estable
(0.12~$\rightarrow$~0.10), pero su promedio diario empeor\'o (0.15~$\rightarrow$~0.34): al
activarse el cobro solo en pico, el uso del Metro se concentra en esas franjas y cae casi
a cero en valle, lo que eleva la desigualdad modal medida sobre el d\'ia completo. A
diferencia de corridas previas, el GestorAMT ejerci\'o plenamente su margen din\'amico:
escal\'o la tarifa base del auto desde \$2.00 hasta el techo de \$3.00 (factor 1.5) en
ambas franjas pico ante densidad alta, encareciendo proporcionalmente a SUV y Carga hasta
\$4.50. El escenario EB acumul\'o una recaudaci\'on de 12\,194\,USD ($\sim$762.15\,USD/h
simulada), desglosada por tipo de veh\'iculo en la Tabla~\ref{tab:recaudacion}. Las
pruebas $t$ de Student arrojaron significancia \'unicamente para la transferencia modal
($p=0.004$), mientras que el flujo ($p=0.785$) y la velocidad ($p=0.740$) no resultaron
significativos. La Tabla~\ref{tab:londres} contrasta estos hallazgos con las metas del
\textit{benchmark}.

% ---------- FIGURA 7: benchmark Londres (una columna) ----------
\begin{figure}[!t]
  \centering
  \includegraphics[width=\columnwidth]{images/fig7_benchmark_londres.png}
  \caption{Comparaci\'on de los indicadores del modelo EB frente a las metas del London
           Congestion Charge.}
  \label{fig:benchmark}
\end{figure}

% ---------- TABLA 2: validacion vs Londres (ancho completo) ----------
\begin{table*}[!t]
  \centering
  \caption{Validaci\'on contrafactual frente al London Congestion Charge.}
  \label{tab:londres}
  \small
  \renewcommand{\arraystretch}{1.3}
  \begin{tabularx}{\textwidth}{@{}lllXc@{}}
    \toprule
    \textbf{Indicador} & \textbf{Londres 2003} & \textbf{Meta Quito} &
    \textbf{Modelo EB} & \textbf{\textquestiondown Cumple?}\\
    \midrule
    Reducci\'on vehicular (hora pico) & $-27$\,\% & $\geq -20$\,\% & $-17.2$\,\% & \no\\
    Mejora de velocidad (hora pico)   & $+20$\,\% & $>0$\,\%       & $+2.3$\,\%  & \yes\\
    Modal shift a transporte p\'ublico & $\sim$12\,\% & $>5$\,\%   & $\sim$14.3\,\% & \yes\\
    Tarifa pico & \pounds5 $\approx$ \$6.0 & \$2.00--\$3.00 propuesto &
      Auto \$2--3; SUV/Carga \$3--4.5 (din\'amica); Moto/Bus exentos & ---\\
    \bottomrule
  \end{tabularx}
\end{table*}

% ---------- FIGURA 8: equidad y recaudacion (una columna) ----------
\begin{figure}[!t]
  \centering
  \includegraphics[width=\columnwidth]{images/fig8_equidad_recaudacion.png}
  \caption{\'Indice de Gini modal y recaudaci\'on acumulada del escenario EB.}
  \label{fig:equidad}
\end{figure}

% ---------- FIGURA 9: recaudacion por tipo (una columna) ----------
\begin{figure}[!t]
  \centering
  \includegraphics[width=\columnwidth]{images/fig9_tarifa_por_tipo.png}
  \caption{Recaudaci\'on acumulada y n\'umero de pagos por tipo de veh\'iculo (Auto, SUV,
           Carga; Moto y Bus exonerados).}
  \label{fig:tarifatipo}
\end{figure}

% ---------- TABLA 3: recaudacion por tipo (una columna) ----------
\begin{table}[!t]
  \centering
  \caption{Recaudaci\'on acumulada por tipo de veh\'iculo (escenario EB, d\'ia completo).}
  \label{tab:recaudacion}
  \small
  \renewcommand{\arraystretch}{1.25}
  \begin{tabularx}{\columnwidth}{@{}Xrrr@{}}
    \toprule
    \textbf{Tipo} & \textbf{Recaud.\ (USD)} & \textbf{Pagos} & \textbf{Tarifa ef.\ (USD)}\\
    \midrule
    Auto        & 7334.00  & 3378 & 2.17\\
    SUV         & 2583.75  & 1115 & 2.32\\
    Carga       & 2276.63  & 1048 & 2.17\\
    Moto / Bus  & 0.00     & 0    & 0.00 (exon.)\\
    \midrule
    \textbf{Total} & \textbf{12\,194.38} & \textbf{5541} & ---\\
    \bottomrule
  \end{tabularx}
\end{table}

La tarifa media efectiva por pago es similar entre tipos pese a las bases diferenciadas
(SUV/Carga $=1.5\times$ el auto): cuando el factor din\'amico elev\'o las tarifas de SUV y
Carga por encima de la disposici\'on a pagar de muchos de esos conductores, estos se
desviaron o cambiaron de modo en lugar de pagar, de modo que los pagos efectivamente
registrados se concentran en los tramos de tarifa m\'as baja.

% =====================================================================
%  5. DISCUSION   (tiempo presente/pasado)
% =====================================================================
\section{Discusi\'on}
Los patrones emergentes de la simulaci\'on evidenciaron una correspondencia direccional
con las teor\'ias de gesti\'on de demanda de transporte: al imponerse la penalizaci\'on
financiera en EB, los conductores modificaron sus intenciones en funci\'on de su
elasticidad econ\'omica, desvi\'andose hacia la periferia o transfiri\'endose al Metro para
preservar su utilidad. Este comportamiento confirm\'o que el mecanismo BDI oper\'o
correctamente y que el modelo constituye una herramienta de evaluaci\'on v\'alida, con
independencia del signo del resultado de pol\'itica.

No obstante, la magnitud del efecto, aunque relevante, se qued\'o por debajo del referente
de validaci\'on. La evidencia muestra que un peaje urbano reduce el volumen de tr\'afico
que ingresa a la zona central entre un 15 y un 27\,\% y disminuye la congesti\'on en torno
a un 30\,\% durante el primer a\~no, incrementando las velocidades de circulaci\'on
\cite{london2020,london2024,atmos2012}. En el presente modelo, la reducci\'on de flujo en
hora pico ($-17.2$\,\%) se ubic\'o dentro del rango inferior emp\'irico de Londres
($-15$\,\%) pero no alcanz\'o el umbral de validaci\'on del 20\,\%, por lo que \textbf{la
hip\'otesis principal no se confirm\'o}. Tres factores explican esta brecha. Primero, pese
a que el GestorAMT escal\'o la tarifa hasta su techo (\$3.00 para el auto, \$4.50 para
SUV/Carga), las tarifas evaluadas siguieron siendo menores que la de Londres
(\pounds5~$\approx$~\$6.0 USD en 2003); adem\'as, la exoneraci\'on de motos y buses ---los
modos m\'as numerosos tras el auto--- acot\'o la fracci\'on de flota efectivamente
disuadida, traduci\'endose en una elasticidad de demanda d\'ebil. Segundo, se observ\'o un
fuerte efecto de desplazamiento perif\'erico ($+52.7$\,\%), coherente con la literatura que
advierte que parte del tr\'afico se reasigna a v\'ias no cobradas en lugar de desaparecer
\cite{london2024}; este desv\'io explica buena parte de la ca\'ida de flujo interno.
Tercero, el peaje no agrav\'o la inequidad modal dentro de la franja de cobro (Gini pico
0.12~$\rightarrow$~0.10), pero s\'i elev\'o la desigualdad modal medida sobre el d\'ia
completo (Gini promedio 0.15~$\rightarrow$~0.34), al concentrar el uso del Metro en las
horas pico y dejarlo casi nulo en valle; este patr\'on temporal, m\'as que un sesgo
socioecon\'omico directo, es el origen del deterioro del indicador diario.

Se identifica, adem\'as, una diferencia anal\'itica crucial respecto a Londres: la capital
brit\'anica contaba con una red de transporte p\'ublico madura antes del cobro, mientras
que el Metro de Quito opera comercialmente solo desde diciembre de 2023 \cite{metro2025}.
Los datos de calibraci\'on indican que $\sim$15\,\% de los usuarios actuales del Metro
migr\'o desde el veh\'iculo particular ($\sim$10\,000 autos/d\'ia ya removidos), de modo
que la simulaci\'on midi\'o el efecto \textit{adicional} del peaje superpuesto a una
absorci\'on modal base a\'un en consolidaci\'on. Esto sugiere que el bajo rendimiento
observado responde al dise\~no tarifario y a la madurez del sistema de absorci\'on, m\'as
que a una inviabilidad estructural del instrumento.

Como limitaci\'on metodol\'ogica, las velocidades medias simuladas ($\sim$31\,km/h)
excedieron el rango emp\'irico de hora pico del DMQ (14--18\,km/h), por lo que los valores
absolutos deben leerse como tendencias relativas entre escenarios y no como predicciones
puntuales. Asimismo, la colocaci\'on interactiva de los puntos de control en los modelos
configurables limita por ahora la reproducibilidad en ejecuci\'on desatendida.

% =====================================================================
%  6. CONCLUSIONES   (tiempo pasado, ligadas a los objetivos)
% =====================================================================
\FloatBarrier
\section{Conclusiones}
El desarrollo del modelo multiagente integrando GAMA Platform y QGIS~3.x permiti\'o
evaluar con rigor el impacto de una pol\'itica de tarificaci\'on vial en un entorno urbano
andino. Respecto al \textbf{primer objetivo}, se concluy\'o que la ingesta de las redes
viales vectoriales reales desde OpenStreetMap constituy\'o una plataforma geom\'etrica de
alta fidelidad, adecuada para el monitoreo de los flujos continuos. En relaci\'on con el
\textbf{segundo objetivo}, los perfiles basados en la arquitectura BDI capturaron
exitosamente las respuestas conductuales heterog\'eneas ante est\'imulos econ\'omicos
disuasorios, superando las limitaciones de los enfoques promediados; el mecanismo
deliberativo oper\'o de forma consistente en ambos escenarios. Atendiendo al \textbf{tercer
objetivo}, la comparaci\'on sistem\'atica determin\'o que la zona de cobro disminuy\'o el
tr\'afico de paso interno en hora pico ($-17.2$\,\%) e increment\'o significativamente la
demanda del Metro ($+45$\,\%, $p=0.004$), con el agente gestor ejerciendo plenamente su
rango din\'amico hasta el techo tarifario; pero \textbf{no alcanz\'o el umbral de
reducci\'on del 20\,\%} y gener\'o una externalidad de desplazamiento perif\'erico
($+52.7$\,\%) y una concentraci\'on temporal del uso del Metro que deterior\'o la equidad
modal en el promedio diario.

En consecuencia, bajo los par\'ametros evaluados, el modelo no respald\'o la hip\'otesis
principal: el peaje propuesto produjo beneficios moderados con efectos colaterales
relevantes. Este hallazgo, lejos de invalidar el instrumento, demostr\'o su utilidad como
herramienta predictiva de soporte a la decisi\'on para el Plan Maestro de Movilidad del
Municipio de Quito, capaz de anticipar que la medida ---tal como est\'a dise\~nada---
rendir\'ia poco. Como \textbf{trabajo futuro} se plante\'o: (i)~elevar el techo tarifario
del gestor por encima de los \$3.00 actuales hacia el rango \$3--6 USD y revisar la
exoneraci\'on de motos ---que concentran parte del flujo disuasible--- reejecutando EB;
(ii)~revisar la distribuci\'on de disposici\'on a pagar y los umbrales de congesti\'on por
NSE; (iii)~incorporar mecanismos de mitigaci\'on de equidad (exenciones por estrato
socioecon\'omico); (iv)~calibrar el modelo con aforos reales y externalidades ambientales;
y (v)~extender el an\'alisis a los escenarios de peaje fijo 24\,h y tarifa din\'amica en
tiempo real.

% =====================================================================
%  BIBLIOGRAFIA  (BibTeX: ver references.bib)
%  Compilar con: pdflatex main -> bibtex main -> pdflatex main -> pdflatex main
% =====================================================================
\bibliographystyle{ieeetr}
\small
\bibliography{references}

\end{document}
